\section{Quasi-subsequentialization of functional separated finite-state transducers}
\label{section:quasi_subsequentialization}

We now describe an inductive construction that, given a functional $k$-separated
finite-state transducer $\mathcal{T} = \langle \Sigma, \Omega, Q, X, S, F, \Delta
\rangle$, produces a $k$-separated quasi-subsequential finite-state transducer
$\mathcal{T}'$ such that $\mathcal{R}_{\mathcal{T}'} = \mathcal{F}_\mathcal{T}$.

The quasi-states of $\mathcal{T}'$ are sets of pairs $\langle p, w \rangle \in Q
\times \Omega^*$ of a state and a pending output. For a set $P \subseteq Q \times
\Omega^*$ we write $out(P) = \{ w \in \Omega^* \mid \exists p \in Q (\langle p, w
\rangle \in P) \}$. We extend the longest common prefix by a symbol $\bot$ for the
empty set: $\bigwedge \emptyset = \bot$, $u \wedge \bot = \bot \wedge u = u$ for $u \in
\Omega^* \cup \{ \bot \}$, and $\bot^{-1}w = w$ for $w \in \Omega^*$. For $m \in \Omega^*
\cup \{ \bot \}$ and $P \subseteq Q \times \Omega^*$ such that $m$ is a prefix of every
word in $out(P)$, we write $m^{-1}P = \{ \langle p, m^{-1}w \rangle \mid \langle p, w
\rangle \in P \}$.

The separator states of $\mathcal{T}'$ are pairs $\langle m, E \rangle \in X'$ where
\begin{equation*}
  X' = (\Omega^* \cup \{ \bot \}) \times 2^{X \times \Omega^*}.
\end{equation*}
The first component carries the longest common prefix of the outputs of a
quasi-state, and the second one its pairs with separator states. Passing the longest
common prefix through the separator states allows every component to normalize its
pending outputs by the common output $\lambda(q, \sigma)$ of all components, although
each table of $\mathcal{T}'$ only depends on its own arguments.

The start state $s$ of the $k$-separated quasi-subsequential finite-state transducer
$\mathcal{T}'$ is defined as
\begin{equation*}
  \langle s^1, s^2, \ldots, s^k \rangle
\end{equation*}
where $s^i = \hat{\Delta}_\epsilon^i(S \times \{ \epsilon \})$ for $i \in \{ 1, 2,
\ldots, k \}$. The base $k$-separated quasi-subsequential finite-state transducer
$\mathcal{T}'^{(1)}$ is defined as
\begin{equation*}
  \langle \Sigma, \Omega, Q'^{(1)}, X', s, F'^{(1)}, \tau^{(1)}, \rho^{(1)},
  \eta^{(1)}, \emptyset, \lambda^{(1)}, \psi^{(1)} \rangle
\end{equation*}
where $Q_i'^{(1)} = \{ s^i \}$ and $F_i'^{(1)} = \tau_i^{(1)} = \rho_i^{(1)} =
\eta_i^{(1)} = \lambda_i^{(1)} = \psi_i^{(1)} = \emptyset$ for $i \in \{ 1, 2,
\ldots, k \}$.
The base set of states $R^{(1)}$ is defined as $\{ s \}$.
Let us assume that $\mathcal{T}'^{(n)}$ and $R^{(n)}$ are constructed for $n > 0$.
The set of states introduced at step $n$ is $R_n = R^{(n)} \setminus R^{(n-1)}$ where
$R^{(0)} = \emptyset$.
The next quasi-state transition function $\tau_i^{(n+1)}$ is defined as
\begin{equation*}
  \tau_i^{(n)} \cup \{ \langle q^i, \sigma, \Delta_{f\epsilon}^i(q^i, \sigma) \rangle
  \mid \exists q \in R_n \; \exists \sigma \in \Sigma \}
\end{equation*}
for $i \in \{ 1, 2, \ldots, k \}$.
The next quasi-state transition output function $\lambda_i^{(n+1)}$ is defined as
\begin{equation*}
  \lambda_i^{(n)} \cup \{ \langle q^i, \sigma, \bigwedge
  out(\Delta_{f\epsilon}^i(q^i, \sigma)) \rangle \mid \exists q \in R_n \; \exists
  \sigma \in \Sigma (\Delta_{f\epsilon}^i(q^i, \sigma) \neq \emptyset) \}
\end{equation*}
for $i \in \{ 1, 2, \ldots, k \}$.
The next quasi-state reduction function $\rho_i^{(n+1)}$ is defined as
\begin{equation*}
  \rho_i^{(n)} \cup \{ \langle t, \langle \bigwedge out(t), t \cap (X \times \Omega^*)
  \rangle \rangle \mid \exists t \in rng(\tau_i^{(n+1)}) \setminus rng(\tau_i^{(n)}) \}
\end{equation*}
for $i \in \{ 1, 2, \ldots, k \}$.
The next separator extension function $\eta_x^{(n+1)}$ is defined as
\begin{equation*}
  \eta_x^{(n)} \cup \bigcup_{i=1}^{k-1} \{ \langle y_i(q, \sigma), r_{i+1}(q,
  \sigma), y_{i+1}(q, \sigma) \rangle \mid \exists q \in R_n \; \exists \sigma \in
  \Sigma \}
\end{equation*}
where $r_j(q, \sigma) = \rho_j^{(n+1)}(\tau_j^{(n+1)}(q^j, \sigma)) = \langle
m_j, E_j \rangle$ and $y_j(q, \sigma) = \langle m_1 \wedge \ldots \wedge m_j, E_1
\cup \ldots \cup E_j \rangle$ for $j \in \{ 1, 2, \ldots, k \}$. That is, $\eta_x$
maps $\langle \langle m, E \rangle, \langle m', E' \rangle \rangle$ to $\langle m
\wedge m', E \cup E' \rangle$.
The next quasi-state extension function $\eta_i^{(n+1)}$ is defined as
\begin{equation*}
  \eta_i^{(n)} \cup \{ \langle t_i(q, \sigma), x(q, \sigma), m^{-1}(t_i(q, \sigma)
  \cup \hat{\Delta}_\epsilon^i(E)) \rangle \mid \exists q \in R_n \; \exists \sigma
  \in \Sigma \}
\end{equation*}
where $t_i(q, \sigma) = \tau_i^{(n+1)}(q^i, \sigma)$ and $x(q, \sigma) =
\delta_x^{(n+1)}(q, \sigma) = \langle m, E \rangle$ for $i \in \{ 1, 2, \ldots, k
\}$.
The next set of final quasi-states $F_i'^{(n+1)}$ is defined as
\begin{equation*}
  F_i'^{(n)} \cup \{ q^i \mid \exists q \in R_n \; \exists \langle p, w \rangle \in
  q^i (p \in F) \}
\end{equation*}
for $i \in \{ 1, 2, \ldots, k \}$.
The next final quasi-state output function $\psi_i^{(n+1)}$ extends $\psi_i^{(n)}$ by
$\psi_i^{(n+1)}(q^i) = w$ for every $q^i \in F_i'^{(n+1)} \setminus F_i'^{(n)}$, where
$\langle p, w \rangle \in q^i$ with $p \in F$ is chosen arbitrarily. By
\Cref{prop:quasisubseq_psi} below, the choice does not matter.
The next set of states $Q_i'^{(n+1)}$ is defined as
\begin{equation*}
  Q_i'^{(n)} \cup rng(\tau_i^{(n+1)}) \cup rng(\eta_i^{(n+1)})
\end{equation*}
for $i \in \{ 1, 2, \ldots, k \}$.
The next $k$-separated quasi-subsequential finite-state transducer
$\mathcal{T}'^{(n+1)}$ is defined as
\begin{equation*}
  \langle \Sigma, \Omega, Q'^{(n+1)}, X', s, F'^{(n+1)}, \tau^{(n+1)},
  \rho^{(n+1)}, \eta^{(n+1)}, \eta_x^{(n+1)}, \lambda^{(n+1)}, \psi^{(n+1)} \rangle.
\end{equation*}
The next set of states $R^{(n+1)}$ is defined as
\begin{equation*}
  R^{(n)} \cup \{ \delta^{(n+1)}(q, \sigma) \mid \exists q \in R_n \; \exists \sigma
  \in \Sigma (\lambda^{(n+1)}(q, \sigma) \textrm{ is defined}) \}.
\end{equation*}

\begin{proposition}
  Let $\mathcal{T} = \langle \Sigma, \Omega, Q, X, S, F, \Delta \rangle$ be a
  $k$-separated finite-state transducer. Let $\mathcal{T}'^{(n)}$ be the
  $k$-separated quasi-subsequential finite-state transducer derived from
  $\mathcal{T}$ by applying the construction described above for $n > 0$. Then
  $\tau_i^{(n)}, \rho_i^{(n)}, \eta_i^{(n)}, \lambda_i^{(n)}$ for $i \in \{ 1, 2,
  \ldots, k \}$ and $\eta_x^{(n)}$ are functions.
\end{proposition}

\begin{proof}
  Every entry added to $\tau_i^{(n)}$ or $\lambda_i^{(n)}$ is determined by its
  arguments $q^i$ and $\sigma$, since $\Delta_{f\epsilon}^i$ is a function. Every entry
  added to $\rho_i^{(n)}$ is determined by its argument $t$. Every entry $\langle
  \langle m, E \rangle, \langle m', E' \rangle, y \rangle$ added to $\eta_x^{(n)}$
  satisfies $y = \langle m \wedge m', E \cup E' \rangle$. Every entry $\langle t,
  \langle m, E \rangle, e \rangle$ added to $\eta_i^{(n)}$ satisfies $e = m^{-1}(t \cup
  \hat{\Delta}_\epsilon^i(E))$. Therefore, all of them are functions.
\end{proof}

The value $\tau_i(q^i, \sigma)$ is not normalized by the common output $\lambda(q,
\sigma)$, because $\lambda(q, \sigma)$ depends on all components of $q$ and not only
on $q^i$. A normalization inside $\tau_i$ would in general not be a function of $q^i$
and $\sigma$. The normalization is instead performed by $\eta_i$, which receives the
common output as the first component of the separator state $\delta_x(q, \sigma)$.

\begin{proposition}
  Let $\mathcal{T} = \langle \Sigma, \Omega, Q, X, S, F, \Delta \rangle$ be a
  $k$-separated finite-state transducer. Let $\mathcal{T}'^{(n)}$ be the
  $k$-separated quasi-subsequential finite-state transducer derived from
  $\mathcal{T}$ by applying the construction described above for $n > 0$. Then
  $dom(\lambda_i^{(n)}) = \{ \langle P, \sigma \rangle \in dom(\tau_i^{(n)}) \mid
  \tau_i^{(n)}(P, \sigma) \neq \emptyset \}$ and $dom(\rho_i^{(n)}) =
  rng(\tau_i^{(n)})$ for $i \in \{ 1, 2, \ldots, k \}$.
\end{proposition}

\begin{proof}
  Both statements hold for $n = 1$, since $\lambda_i^{(1)} = \tau_i^{(1)} =
  \rho_i^{(1)} = \emptyset$. At step $n + 1$, an entry $\langle q^i, \sigma, t
  \rangle$ with $t = \Delta_{f\epsilon}^i(q^i, \sigma)$ is added to $\tau_i$, and an
  entry for $\langle q^i, \sigma \rangle$ is added to $\lambda_i$ if and only if $t
  \neq \emptyset$. Also, $\rho_i$ is extended exactly by the elements of
  $rng(\tau_i^{(n+1)}) \setminus rng(\tau_i^{(n)})$. Thus, both statements follow by
  induction on $n$.
\end{proof}

\begin{proposition}\label{prop:quasisubseq_delta}
  Let $\mathcal{T} = \langle \Sigma, \Omega, Q, X, S, F, \Delta \rangle$ be a
  $k$-separated finite-state transducer. Let $\mathcal{T}'^{(n)}$ be the
  $k$-separated quasi-subsequential finite-state transducer derived from
  $\mathcal{T}$ by applying the construction described above for $n > 1$. Let $q \in
  R^{(n-1)}$ and $\sigma \in \Sigma$. Let
  \begin{equation*}
    W = \bigcup_{j=1}^k out(\Delta_{f\epsilon}^j(q^j, \sigma)), \qquad E =
    \bigcup_{j=1}^k (\Delta_{f\epsilon}^j(q^j, \sigma) \cap (X \times \Omega^*)).
  \end{equation*}
  Then $\lambda(q, \sigma)$ is defined if and only if $W \neq \emptyset$. In that
  case, $\lambda(q, \sigma) = \bigwedge W$, $\delta_x(q, \sigma) = \langle \lambda(q,
  \sigma), E \rangle$, and
  \begin{equation*}
    \delta_i(q, \sigma) = \lambda(q, \sigma)^{-1}(\Delta_{f\epsilon}^i(q^i, \sigma)
    \cup \hat{\Delta}_\epsilon^i(E))
  \end{equation*}
  for $i \in \{ 1, 2, \ldots, k \}$.
\end{proposition}

\begin{proof}
  By construction, $\tau_j(q^j, \sigma) = \Delta_{f\epsilon}^j(q^j, \sigma)$, and
  $\lambda_j(q^j, \sigma)$ is defined if and only if this set is nonempty, in which
  case it equals $\bigwedge out(\Delta_{f\epsilon}^j(q^j, \sigma))$. Since the longest
  common prefix of a union of sets is the longest common prefix of their longest
  common prefixes, $\lambda(q, \sigma) = \bigwedge W$ if $W \neq \emptyset$. Let
  $\rho_j(\tau_j(q^j, \sigma)) = \langle m_j, E_j \rangle$. By induction on $i$,
  $\eta_x^i(\rho_1(\tau_1(q^1, \sigma)), \ldots, \rho_i(\tau_i(q^i, \sigma))) =
  \langle m_1 \wedge \ldots \wedge m_i, E_1 \cup \ldots \cup E_i \rangle$. For $i = k$
  this is $\langle \bigwedge W, E \rangle$, since $m_j = \bigwedge
  out(\Delta_{f\epsilon}^j(q^j, \sigma))$ and $\bot$ is neutral for $\wedge$.
  Finally, $\delta_i(q, \sigma) = \eta_i(\tau_i(q^i, \sigma), \delta_x(q, \sigma))$
  by definition.
\end{proof}

Since the entries added for a state depend only on that state, the tables of
$\mathcal{T}'^{(n)}$ are restrictions of the tables of $\mathcal{T}'^{(n+1)}$. Hence,
for $w \in \Sigma^*$ with $|w| < n$, the values $\delta^{(n)*}(s, w)$ and
$\lambda^{(n)*}(s, w)$, and whether they are defined, do not depend on $n$. We denote
them by $\delta^*(s, w)$ and $\lambda^*(s, w)$. By induction on $n$, $R^{(n)} = \{
\delta^*(s, w) \mid w \in \Sigma^*, |w| < n, \delta^*(s, w) \textrm{ is defined} \}$.
For $w \in \Sigma^*$ let
\begin{equation*}
  Run(w) = \{ \langle p, v \rangle \in Q \times \Omega^* \mid \exists s' \in S
  (\langle s', w, v, p \rangle \in \Delta^*) \}.
\end{equation*}

\begin{lemma}\label{lem:quasisubseq_invariant}
  Let $\mathcal{T} = \langle \Sigma, \Omega, Q, X, S, F, \Delta \rangle$ be a
  $k$-separated finite-state transducer, and let $w \in \Sigma^*$. Then
  \begin{enumerate}
    \item $\lambda^*(s, w)$ is defined if and only if $w = \epsilon$ or $Run(w) \neq
      \emptyset$,
    \item if $o = \lambda^*(s, w)$ is defined, then $o$ is a prefix of every word in
      $out(Run(w))$ and
      \begin{equation*}
        \delta^*(s, w)^i = \{ \langle p, o^{-1}v \rangle \mid \langle p, v \rangle \in
        Run(w), p \in \mathcal{C}(Q_i) \}
      \end{equation*}
      for $i \in \{ 1, 2, \ldots, k \}$,
    \item if $o = \lambda^*(s, w)$ is defined and $w \neq \epsilon$, then $o =
      \bigwedge out(Run(w))$.
  \end{enumerate}
\end{lemma}

\begin{proof}
  We prove the statement by induction on the length of $w$. If $|w| = 0$, then $w =
  \epsilon$, $\lambda^*(s, \epsilon) = \epsilon$ and
  \begin{equation*}
    \delta^*(s, \epsilon)^i = s^i = \hat{\Delta}_\epsilon^i(S \times \{ \epsilon \}) =
    \{ \langle p, v \rangle \mid \langle p, v \rangle \in Run(\epsilon), p \in
    \mathcal{C}(Q_i) \}.
  \end{equation*}

  Otherwise, $w = u\sigma$ for $u \in \Sigma^*$ and $\sigma \in \Sigma$. If
  $\lambda^*(s, u)$ is undefined, then so is $\lambda^*(s, w)$, and by the inductive
  hypothesis $Run(u) = \emptyset$. Every path with label $\langle u\sigma, v
  \rangle$ from an initial state starts with a path with label $\langle u, v'
  \rangle$, so also $Run(w) = \emptyset$.

  Let $\lambda^*(s, u) = o$ be defined, let $p = \delta^*(s, u)$, and let $W$ and $E$
  be as in \Cref{prop:quasisubseq_delta} for $p$ and $\sigma$. By the inductive
  hypothesis, $p^r = \{ \langle q, o^{-1}v \rangle \mid \langle q, v \rangle \in
  Run(u), q \in \mathcal{C}(Q_r) \}$ for $r \in \{ 1, 2, \ldots, k \}$. We make two
  observations.

  First, if $\langle q, d \rangle \in \Delta_{f\epsilon}^r(p^r, \sigma)$ for some $r$,
  then $\langle q, od \rangle \in Run(w)$. Indeed, there exist $\langle q', d' \rangle
  \in p^r$ and $v \in \Omega^*$ such that $\langle q', \sigma, v, q \rangle \in
  \Delta^*$ and $d = d'v$. Since $\langle q', od' \rangle \in Run(u)$, also $\langle
  q, od'v \rangle \in Run(u\sigma)$.

  Second, let $\langle q, v \rangle \in Run(w)$ with $q \in \mathcal{C}(Q_i)$. Then $v
  = od$ for some $d \in \Omega^*$ such that either $\langle q, d \rangle \in
  \Delta_{f\epsilon}^i(p^i, \sigma)$, or there exist $\langle x, d_1 \rangle \in E$
  and $d_2 \in \Omega^*$ such that $d = d_1 d_2$ and $\langle x, \epsilon, d_2, q
  \rangle \in \Delta^*$. Indeed, consider a path with label $\langle w, v \rangle$
  from an initial state $s' \in S$ to $q$, and the last transition on it with input
  $\sigma$. There exist $q', q'' \in Q$ and $v_1, v_2, v_3 \in \Omega^*$ such that $v
  = v_1 v_2 v_3$, $\langle s', u, v_1, q' \rangle \in \Delta^*$, $\langle q', \sigma,
  v_2, q'' \rangle \in \Delta$ and $\langle q'', \epsilon, v_3, q \rangle \in
  \Delta^*$. Since $\langle q', v_1 \rangle \in Run(u)$, by the inductive hypothesis
  $v_1 = od'$ for some $d' \in \Omega^*$. If $q' \in \mathcal{C}(Q_i)$, then $\langle
  q', d' \rangle \in p^i$ and $\langle q, d'v_2v_3 \rangle \in
  \Delta_{f\epsilon}^i(p^i, \sigma)$. Otherwise, $q' \in Q_r$ for some $r \neq i$.
  Since $\mathcal{T}$ is $k$-separated, $q'' \in \mathcal{C}(Q_r)$. By
  \Cref{prop:fst_separator_path} applied to $\langle q'', \epsilon, v_3, q \rangle$,
  there exist $x \in X$ and $v_3', v_3'' \in \Omega^*$ such that $v_3 = v_3'v_3''$,
  $\langle q'', \epsilon, v_3', x \rangle \in \Delta^*$ and $\langle x, \epsilon,
  v_3'', q \rangle \in \Delta^*$. Since $\langle q', d' \rangle \in p^r$ and $x \in X
  \subseteq \mathcal{C}(Q_r)$, we have $\langle x, d'v_2v_3' \rangle \in
  \Delta_{f\epsilon}^r(p^r, \sigma) \cap (X \times \Omega^*) \subseteq E$. The
  statement follows with $d_1 = d'v_2v_3'$ and $d_2 = v_3''$.

  By the first observation, $W \neq \emptyset$ implies $Run(w) \neq \emptyset$. Since
  every state belongs to some component $\mathcal{C}(Q_i)$, by the second observation
  $Run(w) \neq \emptyset$ implies $W \neq \emptyset$. By
  \Cref{prop:quasisubseq_delta}, $\lambda(p, \sigma)$ is defined if and only if $W
  \neq \emptyset$. This proves the first statement.

  Let $W \neq \emptyset$ and $m = \lambda(p, \sigma) = \bigwedge W$. Then
  $\lambda^*(s, w) = om$. Let $\langle q, v \rangle \in Run(w)$ and let $d$ be as in
  the second observation. Since $d \in W$, or $d = d_1d_2$ with $d_1 \in W$, $m$ is a
  prefix of $d$, and thus $om$ is a prefix of $v$. By \Cref{prop:quasisubseq_delta},
  \begin{equation*}
    \delta^*(s, w)^i = \delta_i(p, \sigma) = m^{-1}(\Delta_{f\epsilon}^i(p^i, \sigma)
    \cup \hat{\Delta}_\epsilon^i(E)).
  \end{equation*}
  Let $\langle q, e \rangle \in \delta^*(s, w)^i$. Then $q \in \mathcal{C}(Q_i)$ and
  $e = m^{-1}d$, where either $\langle q, d \rangle \in \Delta_{f\epsilon}^i(p^i,
  \sigma)$, or $d = d_1 d_2$ for some $\langle x, d_1 \rangle \in E$ and $\langle x,
  \epsilon, d_2, q \rangle \in \Delta^*$. In the first case, $\langle q, od \rangle
  \in Run(w)$ by the first observation. In the second case, $\langle x, od_1 \rangle
  \in Run(w)$ by the first observation, and appending the path with label $\langle
  \epsilon, d_2 \rangle$ from $x$ to $q$ yields $\langle q, od_1d_2 \rangle = \langle
  q, od \rangle \in Run(w)$. In both cases, $m$ is a prefix of $d$, and $e = m^{-1}d =
  (om)^{-1}(od)$ by \Cref{proposition:left_word_factor}. Conversely, let $\langle q, v \rangle \in
  Run(w)$ with $q \in \mathcal{C}(Q_i)$, and let $d$ be as in the second observation.
  Then $\langle q, d \rangle \in \Delta_{f\epsilon}^i(p^i, \sigma) \cup
  \hat{\Delta}_\epsilon^i(E)$, and thus $\langle q, m^{-1}d \rangle = \langle q,
  (om)^{-1}v \rangle \in \delta^*(s, w)^i$. This proves the second statement.

  Finally, $om$ is a common prefix of $out(Run(w))$. Let $c \in \Omega^*$ be such that
  $omc$ is also a common prefix of $out(Run(w))$. For every $d \in W$, $od \in
  out(Run(w))$ by the first observation, and thus $mc$ is a prefix of $d$. Since $m =
  \bigwedge W$, it follows that $c = \epsilon$. Since any two common prefixes of
  $out(Run(w))$ are prefixes of the same word, one of them is a prefix of the other.
  Therefore, $om = \bigwedge out(Run(w))$, which proves the third statement.
\end{proof}

\begin{proposition}\label{prop:quasisubseq_psi}
  Let $\mathcal{T} = \langle \Sigma, \Omega, Q, X, S, F, \Delta \rangle$ be a
  functional $k$-separated finite-state transducer. Let $q = \delta^*(s, u)$ for some
  $u \in \Sigma^*$, and let $\langle p, w \rangle, \langle p', w' \rangle \in q^i$ with
  $p, p' \in F$ for some $i \in \{ 1, 2, \ldots, k \}$. Then $w = w'$. In particular,
  $\psi_i^{(n)}$ does not depend on the choices made in the construction.
\end{proposition}

\begin{proof}
  Let $o = \lambda^*(s, u)$. By \Cref{lem:quasisubseq_invariant}, $\langle p, ow
  \rangle, \langle p', ow' \rangle \in Run(u)$. Thus, $\langle u, ow \rangle, \langle
  u, ow' \rangle \in \mathcal{R}_\mathcal{T}$. Since $\mathcal{T}$ is functional, $ow
  = ow'$, and thus $w = w'$.
\end{proof}

\begin{lemma}\label{lem:quasisubseq_completion}
  Let $\mathcal{T} = \langle \Sigma, \Omega, Q, X, S, F, \Delta \rangle$ be a trimmed
  functional $k$-separated finite-state transducer such that
  $\mathcal{F}_\mathcal{T}$ has bounded variation. Let $\mathcal{T'}^{(n)}$ and
  $R^{(n)}$ be respectively the $k$-separated quasi-subsequential finite-state
  transducer and the set of states derived from $\mathcal{T}$ by applying the
  construction described above for $n > 0$. There exists $m \in \mathbb{N}$ such that
  $R^{(n)} = R^{(n+1)}$ and $\mathcal{T}'^{(n)} = \mathcal{T}'^{(n+1)}$ for $n > m$.
\end{lemma}

\begin{proof}
  Since $\mathcal{T}$ is trimmed, we can assign each state $q \in Q$ a path $\langle
  q, w_q, v_q, f_q \rangle \in \Delta^*$ with $f_q \in F$. Let $W = \max_{q \in Q}
  |w_q|$ and $V = \max_{q \in Q} |v_q|$. Since $\mathcal{F}_\mathcal{T}$ has bounded
  variation, for $c = 2W$ there exists $K \geq 0$ such that $d_\Sigma(w', w'') \leq
  c$ implies $d_\Omega(\mathcal{F}_\mathcal{T}(w'), \mathcal{F}_\mathcal{T}(w'')) \leq
  K$ for $w', w'' \in dom(\mathcal{F}_\mathcal{T})$.

  We show that $|d| \leq K + V$ for every $w \in \Sigma^* \setminus \{ \epsilon \}$
  such that $\delta^*(s, w)$ is defined and every $\langle p, d \rangle \in
  \delta^*(s, w)^i$, $i \in \{ 1, 2, \ldots, k \}$. Let $o = \lambda^*(s, w)$ and let
  us assume that $|d| > K + V$. Then $d = ad'$ for some $a \in \Omega$ and $d' \in
  \Omega^*$. By \Cref{lem:quasisubseq_invariant}, $\langle p, od \rangle \in Run(w)$
  and $o = \bigwedge out(Run(w))$. Thus, $oa$ is not a common prefix of $out(Run(w))$,
  and there exists $\langle p'', v'' \rangle \in Run(w)$ such that $oa$ is not a
  prefix of $v''$. Again by \Cref{lem:quasisubseq_invariant}, $v'' = od''$ for some
  $d'' \in \Omega^*$, and $a$ is not a prefix of $d''$. The pair $\langle p'', d''
  \rangle$ may belong to a different component than $\langle p, d \rangle$. Since
  $\langle p, od \rangle, \langle p'', od'' \rangle \in Run(w)$, we have
  \begin{equation*}
    \mathcal{F}_\mathcal{T}(ww_p) = odv_p, \qquad \mathcal{F}_\mathcal{T}(ww_{p''}) =
    od''v_{p''}.
  \end{equation*}
  Moreover, $d_\Sigma(ww_p, ww_{p''}) \leq |w_p| + |w_{p''}| \leq 2W = c$. Since the
  common prefix $o$ does not change the sequential distance,
  $d_\Omega(odv_p, od''v_{p''}) = d_\Omega(dv_p, d''v_{p''})$. If $d'' = \epsilon$,
  then $|dv_p \wedge v_{p''}| \leq |v_{p''}| \leq V$, and thus
  \begin{equation*}
    d_\Omega(dv_p, v_{p''}) \geq |dv_p| + |v_{p''}| - 2|v_{p''}| \geq |d| - V > K.
  \end{equation*}
  Otherwise, $d$ and $d''$ start with different symbols, $dv_p \wedge d''v_{p''} =
  \epsilon$, and thus
  \begin{equation*}
    d_\Omega(dv_p, d''v_{p''}) = |dv_p| + |d''v_{p''}| \geq |d| > K.
  \end{equation*}
  This contradicts the choice of $K$. Therefore, $|d| \leq K + V$.

  Let $\Omega^{\leq K+V}$ be the finite set of words over $\Omega$ of length at most $K
  + V$. It follows that every state in $R^{(n)} \setminus \{ s \}$ belongs to the
  finite set $\prod_{i=1}^k 2^{\mathcal{C}(Q_i) \times \Omega^{\leq K+V}}$. Since
  $R^{(n)} \subseteq R^{(n+1)}$ for $n > 0$, there exists $m > 0$ such that $R^{(m)} =
  R^{(m+1)}$, and thus $R_{m+1} = \emptyset$. In that case, $\mathcal{T}'^{(m+2)} =
  \mathcal{T}'^{(m+1)}$ and $R^{(m+2)} = R^{(m+1)}$. It follows that $R^{(n)} =
  R^{(n+1)}$ and $\mathcal{T}'^{(n)} = \mathcal{T}'^{(n+1)}$ for $n > m$.
\end{proof}

The proof of \Cref{lem:quasisubseq_completion} uses the functionality of $\mathcal{T}$
only to express the bounded variation of $\mathcal{F}_\mathcal{T}$. The same argument
applies to any trimmed $k$-separated finite-state transducer whose relation
$\mathcal{R}_\mathcal{T}$ satisfies the bounded variation condition for pairs in
$\mathcal{R}_\mathcal{T}$.

For the smallest $n > 0$ such that $\mathcal{T}'^{(n)} = \mathcal{T}'^{(n+1)}$ we
define $\mathcal{T}' = \mathcal{T}'^{(n)}$ and $R = R^{(n)}$. By
\Cref{lem:quasisubseq_invariant}, $R = \{ \delta^*(s, w) \mid w \in \Sigma^*,
\delta^*(s, w) \textrm{ is defined} \}$.

\begin{lemma}\label{lem:quasisubseq_inclusion_right}
  Let $\mathcal{T} = \langle \Sigma, \Omega, Q, X, S, F, \Delta \rangle$ be a trimmed
  functional $k$-separated finite-state transducer such that
  $\mathcal{F}_\mathcal{T}$ has bounded variation.
  Let $\mathcal{T}' = \langle \Sigma, \Omega, Q', X', s, F', \tau, \rho, \eta,
  \eta_x, \lambda, \psi \rangle$ be the $k$-separated quasi-subsequential
  finite-state transducer derived from $\mathcal{T}$ by applying the construction
  described above. Then $\mathcal{F}_\mathcal{T} \subseteq \mathcal{R}_\mathcal{T'}$.
\end{lemma}

\begin{proof}
  Let $\langle u, v \rangle \in \mathcal{F}_\mathcal{T}$. There exist $s' \in S$ and
  $f \in F$ such that $\langle s', u, v, f \rangle \in \Delta^*$. Thus, $\langle f, v
  \rangle \in Run(u)$. By \Cref{lem:quasisubseq_invariant}, $o = \lambda^*(s, u)$ and
  $\delta^*(s, u) \in R$ are defined, and $v = od$ for some $d \in \Omega^*$. Since $Q
  = X \cup \bigcup_{i=1}^k Q_i$, there exists $i \in \{ 1, 2, \ldots, k \}$ such that
  $f \in \mathcal{C}(Q_i)$. By \Cref{lem:quasisubseq_invariant}, $\langle f, d
  \rangle \in \delta^*(s, u)^i$. Thus, by construction, $\delta^*(s, u)^i \in F_i'$.
  By \Cref{prop:quasisubseq_psi}, $\psi_i(\delta^*(s, u)^i) = d$. Therefore,
  \begin{equation*}
    \langle u, \lambda^*(s, u) \psi_i(\delta^*(s, u)^i) \rangle = \langle u, od
    \rangle = \langle u, v \rangle \in \mathcal{R}_{\mathcal{T}'}.
  \end{equation*}
\end{proof}

\begin{lemma}\label{lem:quasisubseq_inclusion_left}
  Let $\mathcal{T} = \langle \Sigma, \Omega, Q, X, S, F, \Delta \rangle$ be a trimmed
  functional $k$-separated finite-state transducer such that
  $\mathcal{F}_\mathcal{T}$ has bounded variation.
  Let $\mathcal{T}' = \langle \Sigma, \Omega, Q', X', s, F', \tau, \rho, \eta,
  \eta_x, \lambda, \psi \rangle$ be the $k$-separated quasi-subsequential
  finite-state transducer derived from $\mathcal{T}$ by applying the construction
  described above. Then $\mathcal{R}_{\mathcal{T}'} \subseteq
  \mathcal{F}_\mathcal{T}$.
\end{lemma}

\begin{proof}
  Let $\langle u, v \rangle \in \mathcal{R}_{\mathcal{T}'}$. There exists $i \in \{ 1,
  2, \ldots, k \}$ such that $\delta^*(s, u)^i \in F_i'$ and $v = \lambda^*(s, u)
  \psi_i(\delta^*(s, u)^i)$. Let $o = \lambda^*(s, u)$. By construction, there exists
  $\langle p, d \rangle \in \delta^*(s, u)^i$ such that $p \in F$ and
  $\psi_i(\delta^*(s, u)^i) = d$. By \Cref{lem:quasisubseq_invariant}, $\langle p, od
  \rangle \in Run(u)$. Thus, there exists $s' \in S$ such that $\langle s', u, od, p
  \rangle \in \Delta^*$, and $\langle u, v \rangle = \langle u, od \rangle \in
  \mathcal{F}_\mathcal{T}$.
\end{proof}

\begin{theorem}
  For every trimmed functional $k$-separated finite-state transducer $\mathcal{T}$
  such that $\mathcal{F}_\mathcal{T}$ has bounded variation there exists a
  $k$-separated quasi-subsequential finite-state transducer $\mathcal{T'}$ such that
  $\mathcal{R}_{\mathcal{T}'} = \mathcal{F}_\mathcal{T}$.
\end{theorem}

\begin{proof}
  Follows from \Cref{lem:quasisubseq_completion,lem:quasisubseq_inclusion_right,lem:quasisubseq_inclusion_left}.
\end{proof}
